%% ma: L12-B05-0032
%% lop: 12
%% chuong: 1
%% bai: 5
%% dang: 12.5.2
%% loai: TLN
%% muc: TH
%% dap_an: "2"
%% nguon: "HK1-DeOn_1_1.pdf (M10, Đề ôn tập học kì 1 Toán 12, Đề số 2), Phần III Câu 3"
%% kien_thuc: "Bài 5 (ứng dụng đạo hàm giải quyết vấn đề thực tiễn), Lớp 11 Bài 31 (ý nghĩa vật lí của đạo hàm)"
%% trang_thai: chinh-thuc
\begin{ex}
Một chất điểm chuyển động theo phương trình $s=f(t)=-t^3+6t^2+7t+8$, trong đó $s$ tính bằng mét và $t$ tính bằng giây. Thời gian (tính theo giây) để vận tốc của chất điểm đạt giá trị lớn nhất là bao nhiêu?
\shortans{2}
\loigiai{$v(t)=s'(t)=-3t^2+12t+7$, $v'(t)=-6t+12=0\Leftrightarrow t=2$; $v$ đạt giá trị lớn nhất tại $t=2$ (giây).}
\end{ex}
